\subsection*{I. Phương trình trạng thái của khí lí tưởng và phương trình Clapeyron}
immini[thm]{
\newghichu{
$\bullet$ Phương trình trạng thái của một khối lượng khí lí tưởng xác định:
\begin{align*}
	\boder{\dfrac{p_1V_1}{T_1}=\dfrac{p_2V_2}{T_2}\Rightarrow \dfrac{pV}{T}=\text{hằng số}}
\end{align*}
$\bullet$ Độ lớn của hằng số trên phụ thuộc vào lượng khí mà ta xét.
}}
{includegraphics[width=4cm]{Images/KN111T}}
immini[thm]{\begin{itemize}
	\item $n=\dfrac{m}{M}$ với m là khối lượng khí, M là khối lượng mol của lượng khí.
	\item Xét $n=1\ mol$ khí ở điều kiện tiêu chuẩn thì đối với bất kì chất nào cũng đều có thể tích $V=22,4.10^{-3}\ m^3$; áp suất $p=1,013.10^5\ Pa$ và nhiệt độ $T=273\ K$. Vậy:
	\begin{align*}
	\dfrac{pV}{T}=\dfrac{1,013.10^5.22,4.10^{-3}}{273}=8,31	
	\end{align*}
	\item $R=8,31$ được gọi là hằng số khí lí tưởng và có đơn vị là $J/mol.K$.
	\item Vì thể tích của n mol khí bằng n thể tích của 1 mol khí ở cùng nhiệt độ và áp suất nên nếu phương trình trạng thái của 1 mol khí là $\dfrac{pV}{T}=R$ thì phương trình trạng thái của n mol khí là: 
	 \begin{align*}
	 \dfrac{pV}{T}=nR
	 \end{align*}
	\item Phương trình Clapeyron
	\begin{align*}
		\boder{pV=nRT}
	\end{align*}
\end{itemize}}
{includegraphics[width=4cm]{Images/clap}}



\subsection*{II. Vận dụng}
\begin{itemize}
	\item Phương trình trạng thái của khí lí tưởng có nhiều ứng dụng thực tế. 
	\item Nó được dùng vào việc nghiên cứu, chế tạo các thiết bị có liên quan đến chất khí như khí cầu, bình đựng khí, trang phục lặn, máy điều hòa không khí, máy nén khí,...
	\item Nghiên cứu sự thay đổi áp suất và thể tích của các lớp khí tồn tại trong các vật liệu để tìm tòi, sản xuất các vật liệu đáp ứng được các yêu cầu sử dụng khác nhau.
	\item Phương trình này còn được dùng trong việc nghiên cứu sự thay đổi áp suất, nhiệt độ, khối lượng riêng của không khí trong khí quyển, tìm hiểu quá trình biến đổi khí hậu để dự báo thời tiết,...
\end{itemize}